\documentclass[notes,11pt, aspectratio=169]{beamer}

\usepackage{pgfpages}
% These slides also contain speaker notes. You can print just the slides,
% just the notes, or both, depending on the setting below. Comment out the want
% you want.
\setbeameroption{hide notes} % Only slide
%\setbeameroption{show only notes} % Only notes
%\setbeameroption{show notes on second screen=right} % Both

\usepackage{helvet}
\usepackage[default]{lato}
\usepackage{array}
\usepackage{tgbonum}

\usepackage{tikz}
\usepackage{verbatim}
\setbeamertemplate{note page}{\pagecolor{yellow!5}\insertnote}
\usetikzlibrary{positioning}
\usetikzlibrary{snakes}
\usetikzlibrary{calc}
\usetikzlibrary{arrows}
\usetikzlibrary{decorations.markings}
\usetikzlibrary{shapes.misc}
\usetikzlibrary{matrix,shapes,arrows,fit,tikzmark}
\usepackage{amsmath}
\usepackage{mathpazo}
\usepackage{hyperref}
\usepackage{lipsum}
\usepackage{multimedia}
\usepackage{graphicx}
\usepackage{multirow}
\usepackage{graphicx}
\usepackage{dcolumn}
\usepackage{bbm}
\newcolumntype{d}[0]{D{.}{.}{5}}

\usepackage{changepage}
\usepackage{appendixnumberbeamer}
\newcommand{\beginbackup}{
   \newcounter{framenumbervorappendix}
   \setcounter{framenumbervorappendix}{\value{framenumber}}
   \setbeamertemplate{footline}
   {
     \leavevmode%
     \hline
     box{%
       \begin{beamercolorbox}[wd=\paperwidth,ht=2.25ex,dp=1ex,right]{footlinecolor}%
%         \insertframenumber  \hspace*{2ex} 
       \end{beamercolorbox}}%
     \vskip0pt%
   }
 }
\newcommand{\backupend}{
   \addtocounter{framenumbervorappendix}{-\value{framenumber}}
   \addtocounter{framenumber}{\value{framenumbervorappendix}} 
}


\usepackage{graphicx}
\usepackage[space]{grffile}
\usepackage{booktabs}
\newcommand\independent{\protect\mathpalette{\protect\independenT}{\perp}}
\def\independenT#1#2{\mathrel{\rlap{$#1#2$}\mkern2mu{#1#2}}}
\DeclareMathOperator{\Supp}{Supp}


\newtheorem{assN}{Assumption}
% These are my colors -- there are many like them, but these ones are mine.
\definecolor{blue}{RGB}{0,114,178}
\definecolor{red}{RGB}{213,94,0}
\definecolor{yellow}{RGB}{240,228,66}
\definecolor{green}{RGB}{0,158,115}

\hypersetup{
  colorlinks=false,
  linkbordercolor = {white},
  linkcolor = {blue}
}


%% I use a beige off white for my background
\definecolor{MyBackground}{RGB}{255,253,218}

%% Uncomment this if you want to change the background color to something else
%\setbeamercolor{background canvas}{bg=MyBackground}

%% Change the bg color to adjust your transition slide background color!
\newenvironment{transitionframe}{
  \setbeamercolor{background canvas}{bg=yellow}
  \begin{frame}}{
    \end{frame}
}

\setbeamercolor{frametitle}{fg=blue}
\setbeamercolor{title}{fg=black}
\setbeamertemplate{footline}[frame number]
\setbeamertemplate{navigation symbols}{} 
\setbeamertemplate{itemize items}{-}
\setbeamercolor{itemize item}{fg=blue}
\setbeamercolor{itemize subitem}{fg=blue}
\setbeamercolor{enumerate item}{fg=blue}
\setbeamercolor{enumerate subitem}{fg=blue}
\setbeamercolor{button}{bg=MyBackground,fg=blue,}



% If you like road maps, rather than having clutter at the top, have a roadmap show up at the end of each section 
% (and after your introduction)
% Uncomment this is if you want the roadmap!
% \AtBeginSection[]
% {
%    \begin{frame}
%        \frametitle{Roadmap of Talk}
%        \tableofcontents[currentsection]
%    \end{frame}
% }
\setbeamercolor{section in toc}{fg=blue}
\setbeamercolor{subsection in toc}{fg=red}
\setbeamersize{text margin left=1em,text margin right=1em} 

\newenvironment{wideitemize}{\itemize\addtolength{\itemsep}{10pt}}{\enditemize}

\usepackage{environ}
\NewEnviron{videoframe}[1]{
  \begin{frame}
    \vspace{-8pt}
    \begin{columns}[onlytextwidth, T] % align columns
      \begin{column}{.70\textwidth}
        \begin{minipage}[t][\textheight][t]
          {\dimexpr\textwidth}
          \vspace{8pt}
          \hspace{4pt} {\Large \sc \textcolor{blue}{#1}}
          \vspace{8pt}
          
          \BODY
        \end{minipage}
      \end{column}%
      \hfill%
      \begin{column}{.38\textwidth}
        \colorbox{green!20}{\begin{minipage}[t][1.2\textheight][t]
            {\dimexpr\textwidth}
            Face goes here
          \end{minipage}}
      \end{column}%
    \end{columns}
  \end{frame}
}

\title[]{\textcolor{blue}{Canonical Research Designs V:\\ Bartik and Simulated Instruments }}
\author[PGP]{}
\institute[FRBNY]{\small{\begin{tabular}{c}
  Paul Goldsmith-Pinkham  \\
\end{tabular}}}

\date{\today}

\begin{document}

%%% TIKZ STUFF
\tikzset{   
        every picture/.style={remember picture,baseline},
        every node/.style={anchor=base,align=center,outer sep=1.5pt},
        every path/.style={thick},
        }
\newcommand\marktopleft[1]{%
    \tikz[overlay,remember picture] 
        \node (marker-#1-a) at (-.3em,.3em) {};%
}
\newcommand\markbottomright[2]{%
    \tikz[overlay,remember picture] 
        \node (marker-#1-b) at (0em,0em) {};%
}
\tikzstyle{every picture}+=[remember picture] 
\tikzstyle{mybox} =[draw=black, very thick, rectangle, inner sep=10pt, inner ysep=20pt]
\tikzstyle{fancytitle} =[draw=black,fill=red, text=white]
%%%% END TIKZ STUFF

% Title Slide
\begin{frame}
\maketitle
\end{frame}

\begin{frame}{Roadmap for Today}
  \begin{wideitemize}
  \item In some cases, the source of exogeneous variation (either in an IV setting, or just OLS) is straightforward
    \begin{itemize}
    \item There is a single policy or source of variation
    \end{itemize}
  \item However, in other settings, there are more complicated sources of variation exploited to identify effects. Today we'll focus on two:
    \begin{itemize}
    \item Bartik (shift-share) instruments: three recent
      papers on commonly used identification approach
    \item Simulated instruments: reframe an older literature in a new light
    \end{itemize}
  \item Key historical feature of these approaches is that they had an ``intuitive'' feature of identification, but formal properties were not established for several decades
    \begin{itemize}
    \item Analagous to staggered DinD lit!
    \end{itemize}
  \end{wideitemize}
\end{frame}



\begin{frame}{Bartik instruments are used everywhere}
    %Local Markets and National Categories
  \begin{center}
    \begin{tikzpicture}[rectangle,draw=cyan,thick]
      \node (pq1) {Labor supply + immigration};
      \node [dashed,right=of pq1] (pq2) {Trade};
      \node [dashed,right=of pq2] (pq4) {Finance, innovation,\\ others...};
      \node [dashed,above=of pq2] (bartik) {Bartik};
      \path[->,thick,blue] (bartik) edge  (pq1)
      (bartik) edge  (pq2)
      (bartik) edge  (pq4);
    \end{tikzpicture}
  \end{center}
  \begin{itemize}
  \item Thread that links all Bartik applications:
    \begin{itemize}
    \item local markets composed of many ``categories''
    \item need for identification
    \end{itemize}
    \smallskip
  \item Approach has been used since the early 90:
    \begin{itemize}
    \item sometimes called ``shift-share'' or ``industry mix'' instruments
    \end{itemize}
  \end{itemize}
\end{frame}

\begin{frame}{Examples of Bartik instruments in many subfields}
  \begin{description}
  \item[Immigration:] Altonji and Card (1991), Card (2001)
    \smallskip
  \item[Bank Lending:] Amiti and Weinstein (2018), Greenstone, Mas and Nguyen (2015)
    \smallskip    
  \item[Market Size + Demography:] Acemoglu and Linn (2004), Jaravel (2018)
    \smallskip    
  \item[Labor Supply Elasticity:] Blanchard and Katz (1992), \textbf{Bartik} (1991)
    \smallskip
  \item[Fiscal Multipliers:] Nakamura and Steinsson (2014)
    \smallskip    
  \item[Trade + Labor:] Autor, Dorn, and Hanson (2013), Autor, Dorn and Hanson (2018), etc. 
    \smallskip    
  \item[Foreign Aid:] Nunn and Qian (2014)
    \smallskip
  \item[Portfolio Allocation:] Calvet, Campbell, and Sodini (2009)
    \smallskip
  \item[Trade + Prices:] Piveteau and Smagghue (2017), de Roux et al. (2017)
    \smallskip    
  \item[Automation:] Acemoglu and Restrepo (2017)
  \end{description}
\end{frame}

\begin{frame}{Many paths lead to Bartik}
  \begin{itemize}
  \item Diverse literature leads to many motivations and justifications for Bartik approach
    \medskip
  \item Two distinct approaches in the literature:
    \begin{enumerate}
    \item Applied micro statistical approach: interested in a reduced
      form causal relationship; need an instrument that is
      uncorrelated with error term; make argument that Bartik
      instrument is defensible \smallskip
    \item Structural approach: interested in particular parameters
      from model; assumptions of model motivate certain estimating
      equations
    \end{enumerate}
    \medskip
  \item So what is the Bartik approach anyway?
  \end{itemize}
\end{frame}


\begin{frame}
  \frametitle{Motivation: local labor market approaches + reduced form}

  Consider a local labor market regression like the following:  
  \begin{align}
    y_l &= \beta_0 + \beta x_l + \epsilon_l \notag 
  \end{align}

  \begin{itemize}
    \setlength\itemsep{1em}
  \item $\mathbb{E}[x_l\epsilon_l] \neq 0$   $\Rightarrow$ need an instrument to estimate $\beta$ 
  \item E.g. Autor, Dorn and Hanson (2013) setting:
    \begin{itemize}
      \setlength\itemsep{1em}
    \item $l$: location (commuting zone)
    \item $y_l$: manufacturing employment \textit{growth}
    \item  $x_l$: import exposure to China \textit{growth}
    \item $\beta$: effect of rise of China on manufacturing employment
    \item an instrument for location-level exposure to trade with China
    \end{itemize}
\end{itemize}


\end{frame}


%------------
\begin{frame}
\frametitle{The Bartik instrument}

Accounting identity \#1:
\begin{align}
  x_l &=  \sum_{k=1}^{K} z_{lk} g_{lk} \notag  
\end{align}
\begin{itemize}
\setlength\itemsep{1em}
  \item $z_{lk}$: location-industry shares  ($Z_l$)
  \item $g_{lk} $: location-industry growth (in imports) rates ($G_l$) 
  \end{itemize}
Accounting identity \#2:

$$\underbrace{ g_{lk}}_{\text{location-industry}} = \underbrace{  \textcolor{blue}{g}_{ \textcolor{blue}{k}}}_{\text{industry}} + \underbrace{  \textcolor{red}{\tilde{g}_{lk}}   }_{\substack{\text{idiosyncratic} \\ \text{location- industry}}} $$ 

Infeasible Bartik:
$$B_l = \sum_{k=1}^{K}z_{lk} \textcolor{blue}{g_{k}}$$

\end{frame}


%------
\begin{frame}
\frametitle{This gives us a simple 2SLS structure}

\begin{align}
  y_l&= \beta_0 + \beta x_l + \epsilon_l \notag \\
  x_l&= \pi_0 + \pi_{1} B_l + u_l \notag \\  
  B_l &=  \sum_{k=1}^{K}z_{lk} \textcolor{blue}{g_{k}} \notag \\
  g_{lk} &=   \textcolor{blue}{g}_{ \textcolor{blue}{k}}+  \textcolor{red}{\tilde{g}_{lk}}  \notag 
\end{align}



China shock: e.g., Autor, Dorn and Hanson (2013)  
\begin{itemize}
\item $z_{lk}$: location (l) industry (k) composition
\item $g_{lk}$:  location (l)  industry (k) growth in imports from China
\item $\textcolor{blue}{g_k}$: industry $k$  growth of imports (from China)
\end{itemize}


\end{frame}


%%%%%


\begin{frame}
\addtocounter{framenumber}{-1}
\frametitle{Other instruments have this structure}

\begin{align}
  y_l&= \beta_0 + \beta x_l + \epsilon_l \notag \\
  x_l&= \pi_0 + \pi_{1} B_l + u_l \notag \\  
B_l &=  \sum_{k=1}^{K}z_{lk} \textcolor{blue}{g_{k}} \notag \\
g_{lk} &=   \textcolor{blue}{g}_{ \textcolor{blue}{k}}+  \textcolor{red}{\tilde{g}_{lk}}  \notag 
\end{align}

Immigrant enclave: e.g., Altonji and Card (1991)
\begin{itemize}
\item $z_{lk}$: share of people from foreign $k$ living in $l$ (in a base period)
\item $g_{lk}$:  growth in number of people from $k$ to $l$
\item $\textcolor{blue}{g_k}$: growth in people from $k$ nationally
\end{itemize}


\end{frame}


%------
\begin{frame}
\addtocounter{framenumber}{-1}
\frametitle{Other instruments have this structure}

\begin{align}
  y_l&= \beta_0 + \beta x_l + \epsilon_l \notag \\
  x_l&= \pi_0 + \pi_{1} B_l + u_l \notag \\  
B_l &=  \sum_{k=1}^{K}z_{lk} \textcolor{blue}{g_{k}} \notag \\
g_{lk} &=   \textcolor{blue}{g}_{ \textcolor{blue}{k}}+  \textcolor{red}{\tilde{g}_{lk}}  \notag 
\end{align}

Bank-lending relationships: e.g., Greenstone, Mas and Nguyen (2015)
\begin{itemize}
\item $z_{lk}$: location (l) share of loan origination from bank $k$
\item $g_{lk}$:  loan growth in location $l$ by bank $k$
\item $\textcolor{blue}{g_k}$: part of loan growth due to bank supply shock
\end{itemize}


\end{frame}

%------
\begin{frame}
\addtocounter{framenumber}{-1}
\frametitle{Other instruments have this structure}

\begin{align}
  y_l&= \beta_0 + \beta x_l + \epsilon_l \notag \\
  x_l&= \pi_0 + \pi_{1} B_l + u_l \notag \\  
B_l &=  \sum_{k=1}^{K}z_{lk} \textcolor{blue}{g_{k}} \notag \\
g_{lk} &=   \textcolor{blue}{g}_{ \textcolor{blue}{k}}+  \textcolor{red}{\tilde{g}_{lk}}  \notag 
\end{align}

Market size and demography: e.g., Acemoglu and Linn (2004)
\begin{itemize}
\item $z_{lk}$: spending share on drug $l$ from age group $k$
\item $g_{lk}$:  growth in spending of group $k$ on drug $l$
\item $\textcolor{blue}{g_k}$: growth in spending of group $k$ (due to population aging)
\end{itemize}


\end{frame}

\begin{frame}
\frametitle{What's necessary for consistency?}

\begin{align}
  y_l&= \beta_0 + \beta x_l + \epsilon_l \notag \\
  x_l&= \pi_0 + \pi_{1} B_l + u_l \notag \\  
  B_l &=  \sum_{k=1}^{K}z_{lk} \textcolor{blue}{g_{k}} \notag \\
  g_{lk} &=   \textcolor{blue}{g}_{ \textcolor{blue}{k}}+  \textcolor{red}{\tilde{g}_{lk}}  \notag 
\end{align}

\begin{itemize}
\item We need $B_{l}$ to be a valid instrument
\item Requires two conditions with constant effects:
  \begin{enumerate}
  \item Relevance: $\pi_{1} \not=0$, e.g. $Cov(B_{l}, x_{l}) \not= 0$
  \item Exclusion: $E(B_{l}\epsilon_{l}) = 0$
  \end{enumerate}
\item Key flaw in this literature until recently: economic +
  statistical content of exclusion has been vague and sometimes confused
\end{itemize}

\end{frame}

\begin{frame}{Key thing to remember from today}
  \begin{itemize}
  \item Assuming independence or exogeneity on the basis of a model
    does not necessarily make it true
    \begin{itemize}
    \item E.g. Hausman instruments in IO models -- model may assume
      that exclusion restriction is satisfied, but not necessarily
      true in reality
    \end{itemize}
  \item Assuming that two things are independent because they don't
    seem ``related'' doesn't make it true 
    \begin{itemize}
    \item Bartik literature many times argues that national nature of
      shocks ``decouples'' the instrument from local market
      conditions. However, it still exploits local
      characteristics. Need to make very specific arguments to
      validate claim (will come to this).
    \end{itemize}
  \item When evaluating an identification strategy, you should be able
    to describe counterfactual claims using the measure. This is
    typically not concrete in Bartik -- try to make it concrete!
    What is exactly changing in China? Why is it random? 
  \end{itemize}
\end{frame}



%-----
\begin{frame}
\frametitle{More general econometric set-up}
\begin{align}
y_{lt} & =\mathbf{D}_{lt}\beta_0 +x_{lt}\beta+\epsilon_{lt}, \notag \\
x_{lt} & =\mathbf{D}_{lt}\tau+ B_{lt}\gamma+\eta_{lt}\notag \\
\mathbf{D}_{lt} &= \text{\footnotesize{controls, f.e.}} \notag \\
g_{lkt} &=   \textcolor{blue}{g_{ kt}} + \textcolor{red}{\tilde{g}_{lkt}}  \notag \\
B_{lt} &= \sum_{k=1}^K z_{lk0} \textcolor{blue}{g_{kt}}, \notag \\
\{\{  &  x_{lt},\mathbf{D}_{lt},\epsilon_{lt} \}_{t=1}^T\}_{l=1}^L, \text{ iid, } L\rightarrow \infty \notag
\end{align}



  Assumptions for IV in terms of $B_{lt}$:
  \begin{itemize}
	\setlength\itemsep{1em}
  \item Exogeneity: $\mathbb{E}\left[B_{lt}\epsilon_{lt}\middle| \mathbf{D}_{lt}\right] = 0$
  \item Relevance: $\text{Cov}\left[B_{lt}, x_{lt} \middle| \mathbf{D}_{lt}\right] \not= 0$
  \end{itemize}
  
    Question:
  \begin{itemize}
  \item What do these statements about $B_{lt}$ imply about $z_{lk0}$ and $\textcolor{blue}{g_{kt}}$?
%\item Vice-versa?
  \end{itemize}

\end{frame}

\begin{frame}{Recent Literature on this topic}
  \begin{wideitemize}
  \item Three papers addressed this question, and can be split into two grouping
  \item The division between papers can be split based on focus on $z_{lk0}$ vs. $g_{kt}$
    \begin{enumerate}
    \item Goldsmith-Pinkham, Sorkin and Swift (2020) focus on $z_{lk0}$ and make an analogy to difference-in-differences
    \item Adao, Kolesar and Morales (2019) and Borusyak, Hull and Jaravel (2020) focus on $g_{kt}$, and make a strong connection to the design based approach (e.g. these are as-if random shocks)
    \end{enumerate}
  \item I will focus on GPSS, but the important fact is to pick a concrete explanation for identification
    \begin{itemize}
    \item Key problem, historically, in this literature, was the lack
      of a coherent defense of the identifying variation
    \item Ideally, this choice should be made ex ante! I.e. don't run the tests and then choose...
    \end{itemize}
  \end{wideitemize}
\end{frame}


% ----


\begin{frame}
{  \Large
  \begin{center}
    Bartik Instruments: What, When, Why and How\\
    Goldsmith-Pinkham, Sorkin and Swift
  \end{center}
}
  \addtocounter{framenumber}{-1}

Outline:
\begin{itemize}
\setlength\itemsep{1em}
\item \textbf{Understanding the identifying assumption}
\item Opening the black box
\item Testing the plausibility of the identifying assumption
\end{itemize}

\end{frame}


%----
\begin{frame}
\frametitle{Three  special cases}

\begin{enumerate}
  \setlength\itemsep{1em}
\item One time period, two industries
\item T time periods, two industries
\item One time period, K industries
\end{enumerate}

\end{frame}



%------
\begin{frame}[label=twoindustries]
\frametitle{Special case \#1: One time period, two industries}
\begin{itemize}
	\setlength\itemsep{1em}
\item $z_{l2}= 1-z_{l1}$ 
\item Bartik:  \begin{align*}B_l &= z_{l1} \textcolor{blue}{g_1} + z_{l2} \textcolor{blue}{g_2}  =z_{l1} \textcolor{blue}{g_1} + (1-z_{l1}) \textcolor{blue}{g_2} \notag \\
&=\textcolor{blue}{g_2}+ (\textcolor{blue}{g_1} - \textcolor{blue}{g_2}) z_{l1} \notag
\end{align*}
\end{itemize} 
First-stage:
\begin{align*}
x_l &= \gamma_0 + \gamma B_l + \eta_l \notag \\
%x_l &= \gamma_0 + \gamma (\hat{g}_2+ (\hat{g}_1 - \hat{g}_2) z_{l1}) + \eta_l \notag \\
x_l &= \underbrace{\gamma_0+\gamma \textcolor{blue}{g_2}}_{\text{constant}} + \underbrace{\gamma (\textcolor{blue}{g_1} - \textcolor{blue}{g_2})}_{\text{coefficient}} z_{l1} + \eta_l \notag 
\end{align*}

The instrument is $z_{l1}$, while $\textcolor{blue}{g_k}$ affects relevance

\hyperlink{whyolsisbiased}{\beamergotobutton{Why OLS is biased}}




\end{frame}










%---------
\begin{frame}
\frametitle{Special case \#2:  T time periods, two industries }



Panel Bartik:

$$B_{lt} = z_{l10} \textcolor{blue}{g_{1t}}  +z_{l20} \textcolor{blue}{g_{2t}} = \textcolor{blue}{g_{2t}} + \underbrace{ \textcolor{blue}{\Delta_{gt}}}_{g_{1t}-g_{2t}} z_{l10}$$




First stage:
\begin{align}
x_{lt} & =\tau_{l}+\tau_{t}+ \gamma B_{lt} +\eta_{lt} \notag\\
x_{lt} & =\tau_{l}+\underbrace{(\tau_{t}+\gamma \textcolor{blue}{g_{2t}})}_{\tilde{\tau_{t}}}+\underbrace{\gamma  \textcolor{blue}{\Delta_{gt}}}_{\tilde{\gamma}_{t}}z_{l10}+\eta_{lt} \notag
\end{align}




\begin{itemize}
\item Industry shares times time period is the instrument
\item (Updated industry shares: similar)
\end{itemize}


\end{frame}



%---------
\begin{frame}
\frametitle{Special case \#2:  T time periods, two industries }

\begin{itemize}
	\setlength\itemsep{1em}
\item Analogy to continuous difference-in-differences 
	\begin{itemize}
	\setlength\itemsep{1em}
	\item $\textcolor{blue}{\Delta_{gt}}$ is size of policy 
	\item $z_{l10}$ is exposure to policy
	\end{itemize} 
\item Sometimes a ``pre-period'' before policy: test for parallel pre-trends 
	\begin{itemize}
	\setlength\itemsep{1em}
	\item E.g., in ADH, what happens from 1970 to 1990?
	\end{itemize}
\end{itemize}
\end{frame}


%----
\begin{frame}[label=kindustries]
\frametitle{Special case \#3: One time period, K industries}

\begin{itemize}
	\setlength\itemsep{1em}
\item $\textcolor{blue}{G}$: $K \times 1$ vector of $\textcolor{blue}{g_k}$
\item $Z$: $L \times K$, matrix of $Z_l$
\item $Y^\perp$, $X^\perp$, $B = (Z \textcolor{blue}{G})$: $L \times 1$, vectors of $y_l^\perp$, $x_l^\perp$ and $B_l$
\item $\Omega$: $K\times K$
\end{itemize}



\begin{align*}
  \hat{\beta}_{Bartik} &= \frac{B^\prime  Y^\perp}{B^\prime X^\perp}\\
  \hat{\beta}_{GMM} &= \frac{(X^{\perp \prime}Z)\Omega (Z'Y^\perp)}{(X^{\perp \prime} Z)\Omega (Z'X^\perp)}
\end{align*} 

 If $\Omega = (\textcolor{blue}{G} \textcolor{blue}{G}^\prime)$, then $\hat{\beta}_{Bartik}=\hat{\beta}_{GMM}$   

\end{frame}

%-----
\begin{frame}
\frametitle{Full general result with $T$ time periods and $K$ industries}

Two estimators are numerically identical:
\begin{itemize}
\item TSLS with Bartik instrument
\item GMM with industry shares $\times$ time period as instruments and a particular weight matrix 
\end{itemize}
\begin{align*}
  \hat{\beta}_{Bartik} &= \frac{\mathbf{B}^\prime  \widetilde{\mathbf{Y}}^\perp}{\mathbf{B}^\prime \widetilde{\mathbf{X}}^\perp}\\
  \hat{\beta}_{GMM} &= \frac{(\mathbf{X}^{\perp \prime}\widetilde{\mathbf{Z}})\Omega (\widetilde{\mathbf{Z}}'\mathbf{Y}^\perp)}{(\mathbf{X}^{\perp \prime} \widetilde{\mathbf{Z}})\Omega (\widetilde{\mathbf{Z}}'\mathbf{X}^\perp)}
\end{align*} 
$\Omega = (\textcolor{blue}{\mathbf{G}}
\textcolor{blue}{\mathbf{G}}^\prime)$, and $\widetilde{\mathbf{Z}}$ is
an $LT \times KT$ stacked vector of $Z_{0}$ interacted with time fixed
effects and $\textcolor{blue}{\mathbf{G}}$ is a $KT \times 1$ stacked
vector of growth rates $g_{kt}$.

\end{frame}



%----
\begin{frame}
\frametitle{When is the estimator consistent for the estimand of interest?}
\begin{center}
\textbf{What is the identification condition?}
\end{center}

$$\hat{\beta}_{Bartik}=  \frac{  \sum_{l=1}^L \sum_{t=1}^T \sum_{k=1}^K  z_{lkt} \textcolor{blue}{g_{kt}} y_{lt}^\perp    }{   \sum_{l=1}^L \sum_{t=1}^T \sum_{k=1}^K  z_{lkt} \textcolor{blue}{g_{kt}} x_{lt}^\perp   }  $$

Two ideas:
\begin{itemize}
\item ``Shares'' (this paper): talk about properties of $z_{lkt}$
	\begin{itemize}
	\item Conditional exogeneity
	\end{itemize}
\item ``Shocks'' (Borusyak, Hull and Jaravel (2018)): talk about properties of $g_{kt}$
	\begin{itemize}
	\item Random, and a large number (equivalent industry-level regression)
	\end{itemize}
\end{itemize}

\end{frame}

%----
%----
\begin{frame}
\frametitle{When are these views plausible? What do they mean?}


\textbf{Shares}

\textit{Conditional} exogeneity:
\begin{itemize}
\item Typically: exogenous to changes in error term, not levels of outcome
\item Standard in diff-in-diff (exclusion): in a period, exposure to an industry matters for outcome only through $x$
\end{itemize}


\textbf{Shocks}

\begin{itemize}
\item Large number of industries (shares are misspecified, need it to average out)
\item Random shocks across industries
\end{itemize}

\end{frame}

%----
\begin{frame}
\frametitle{Is ADH about shocks or shares?}

Shocks:
\begin{itemize}
\item  Explains why $g_{kt}^{high-income}$ rather than $g_{kt}^{US}$  (hard to rationalize under shares)
\item Natural in a trade model: why would imports from China rise (in a trade model)?  Independent industry-specific shocks
\end{itemize}

Shares:
\begin{itemize}
\item Explains why $z_{lkt-1}$ rather than $z_{lkt}$ (hard to rationalize under shocks)
\item Explains why it is important for identification to study local labor markets (as opposed to parameter of interest where we want to think about spillovers)
\end{itemize}

\medskip

Bottom line: a little hard to tell what exactly ADH are assuming; ADH approach does not appear to satisfy testable assumptions under GPSS, but do appear to under BHJ.
\end{frame}







%%---
\begin{frame}[label=rotemberg]
\frametitle{Decomposing Bartik (GPSS 2020)}
\begin{center}
(Special case of Rotemberg (1983), proposition 1)
\end{center}
\vspace{0.2in}




$$\hat{\beta}_{Bartik} =\sum_k  \hat{\alpha}_k \hat{\beta}_k , \text{      } \sum_k \hat{\alpha}_k=1$$   
IV estimate using only the $k^{th}$ instrument:
$$\hat{\beta}_k = (Z_k^\prime X)^{-1} Z_k^\prime Y$$
``Rotemberg'' weight:
$$ \hat{\alpha}_k = \frac{ \textcolor{blue}{g_k} Z_k^\prime X}{\sum_{k=1}^K \textcolor{blue}{g_k} Z_k^\prime X}$$



\hyperlink{rotproof}{\beamergotobutton{Proof}}

\end{frame}




%----
\begin{frame}
\frametitle{Interpretation: sensitivity to misspecification elasticity}
\begin{center}
  Conley, Hansen and Rossi (2012); Andrews, Gentzkow and Shapiro (2017)\\
  Local misspecification:  $\epsilon_{lt} = L^{-1/2}V_{lt} + \tilde{\epsilon}_{lt}, \enskip Cov(V_{lt}, Z_{lt}) \not=0$,
\end{center}

\begin{itemize}
 \item $   \sqrt{L}\left(\hat{\beta} - \beta_0\right) \xrightarrow{\enskip d \enskip} \tilde{\beta}, \qquad \mathbb{E}\left[\tilde{\beta}\right] = \text{bias (misspecification) of Bartik instrument}$
 \item $   \sqrt{L}\left(\hat{\beta}_{k} - \beta_0\right) \xrightarrow{\enskip d \enskip} \tilde{\beta}_{k}, \qquad \mathbb{E}\left[\tilde{\beta}_{k}\right] = \text{bias (misspecification) of  $k$th  instrument}$

\end{itemize}

Suppose $\beta_0 \neq 0$.  Percentage bias:
\begin{equation}
  \frac{\mathbb{E}[\tilde{\beta}]}{\beta_{0}} = \sum_{k}\alpha_{k}\frac{\mathbb{E}[\tilde{\beta}_{k}]}{\beta_{0}} \notag
\end{equation} 
Industry with high $\alpha_k$: 
\begin{itemize}
\item an industry where it matters whether it is misspecified (endogenous) 
	\begin{itemize}
\item  because it is ``important'' in the estimate
	\end{itemize}
\end{itemize}


\end{frame}



%----
\begin{frame}[label=top5weights]
\textbf{Top five industries (out of 397)}



\begin{center}
\begin{tabular}{lllllllllllll}
\toprule
& $\hat{\alpha}_k$ & $g_k^{\text{high-income}}$ & $\hat{\beta}_k$  \\ \cmidrule(lr){2-4} 
Games and Toys & 0.182 & 174.841 & -0.151  \\  
Electronic Computers & 0.182 & 85.017 & -0.620  \\ 
Household Audio and Video & 0.130 & 118.879 & 0.287  \\ 
Computer Equipment & 0.076 & 28.110 & -0.315  \\ 
Telephone Apparatus & 0.058 & 37.454 & -0.305  \\ 
\bottomrule 
& \multicolumn{2}{l}{$0.628 \big/ 1.379$}   & -0.230 &\\
\end{tabular}
\end{center}
  
\footnotesize
\medskip
\begin{quote}
The \textbf{main source of variation in exposure is} within-manufacturing specialization in industries subject to different degrees of import competition...there is differentiation according to \textbf{local labor market reliance on labor-intensive industries}...By 2007, China accounted for over 40 percent of US imports in four four-digit SIC industries (\textbf{luggage, rubber and plastic footwear, games and toys, and die-cut paperboard}) and over 30 percent in 28 other industries, including \textbf{apparel, textiles, furniture, leather goods, electrical appliances, and jewelry}.
\end{quote}
\begin{center}
--- Autor, Dorn and Hanson (2013) , pg. 2123
\end{center}
\hyperlink{correlmatrix}{\beamergotobutton{Correlation matrix}} 
\end{frame}


%------
\begin{frame}[plain]
\addtocounter{framenumber}{-1}
\frametitle{Outline}


\begin{itemize}
\setlength\itemsep{1em}
\item Understanding the identifying assumption
	\begin{itemize}
	\item \textit{If fixed $K$ and $T$, and} $L \rightarrow \infty$,\textit{ then in terms of industry composition}
	\end{itemize}
\item Opening the black box
	\begin{itemize}
	\item \textit{Which industries are ``important'' in estimates}
	\end{itemize}
\item \textbf{Testing the plausibility of the identifying assumption}
\end{itemize}



\end{frame}


%----
\begin{frame}
\frametitle{Three tests of the identifying condition (under GPSS (2020)}

\begin{enumerate}
\item Confounds (or correlates)
\item Pre-trends 
\item Alternative estimators and overidentification
\end{enumerate}

\begin{wideitemize}
\item There are \emph{also} tests for BHJ -- similar to assuming
  strict ignorability, you can test for balance on observables
  (like the confounds above) of industries and locations
\end{wideitemize}

\end{frame}



%-----
\begin{frame}
\frametitle{Test \#1: Correlates of initial industry composition}

\begin{itemize}
\setlength\itemsep{1em}
\item How are initial characteristics ($D_{l0}$ less F.E.) related to $Z_{l0}$?
\item Look at high-Rotemberg weight industries (and aggregate)
\end{itemize} 

Not definitive:
\begin{itemize}
\setlength\itemsep{1em}
\item Shows source of variation
\item Can address by controlling for observables ($D_{l0} \times \text{time}$)
\end{itemize}

\end{frame}


%-----
\begin{frame}[fragile,label=pca]
\frametitle{Test \#1: Correlates }


\scriptsize

\begin{table}[htbp]
\begin{center}
\begin{tabular}{l*{7}{c}}
\toprule
     &\multicolumn{1}{c}{Games}     &\multicolumn{1}{c}{Electronic}&\multicolumn{1}{c}{Household audio}&\multicolumn{1}{c}{Computer}&\multicolumn{1}{c}{Telephone}&\multicolumn{1}{c}{China }\\
& and toys & computers & and video & equipment & apparatus & to other \\
\midrule
Share Empl in Manufacturing&   0.01&    0.21&     0.08&     0.21&    -0.07&        0.57\\
                &   (0.03)& (0.18)&   (0.08)&   (0.15)&   (0.06)&      (0.07)\\
\marktopleft{a1} Share College Educated&   -0.08&      0.20&     0.01&     0.22&    -0.07&     0.30\\
                &     (0.03)& (0.11)&   (0.04)&   (0.10)&   (0.06)&   (0.06) \markbottomright{a1}{red}  \\
Share Foreign Born& 0.01&   -0.01&    -0.02&    -0.01&    -0.08&          0.15\\
                &  (0.01)& (0.04)&   (0.01)&   (0.04)&   (0.03)&      (0.03)\\
Share Empl of Women& 0.05&       -0.04&    -0.08&    -0.02&    -0.02&      0.10\\
                &  (0.03)&   (0.12)&   (0.05)&   (0.12)&   (0.07)&    (0.06)\\
Share Empl in Routine&0.04&    -0.37&     0.06&    -0.36&    -0.01&         -0.08\\
                &  (0.03)&   (0.14)&   (0.05)&   (0.12)&   (0.07)&    (0.13)\\
Avg Offshorability&    0.02&   0.33&     0.00&     0.29&     0.23&       -0.24\\
                & (0.02)&   (0.10)&   (0.05)&   (0.08)&   (0.04)&     (0.09)\\
\midrule
1980 Population Weighted&      Yes&      Yes&      Yes&      Yes&      Yes&      Yes\\
N               &    1,444&    1,444&    1,444&    1,444&    1,444&    1,444\\
$ R^2$          &    0.02&    0.08&     0.01&     0.08&     0.05&       0.22\\
\bottomrule
\end{tabular}
\end{center}
\end{table}

\uncover<2->{\tikz[overlay,remember picture,inner sep=1pt]
\node[draw=red, line width=0.9mm, rounded corners,fit=(marker-a1-a.north west) (marker-a1-b.south east)] {};}


\end{frame}








%-----

\begin{frame}[label=overidtest]
\frametitle{Test \#2: Pre-trends}


\begin{equation*}
  \Delta \text{Manufacturing Emp}_{lt} = \alpha +  \sum_{s} \mathbb{1}(s=t) \gamma_{k,s} z_{lk, 1980}+ \epsilon_{lt}
\end{equation*}


  \begin{itemize}
\item Four time periods:   1970-1980, 1980-1990, 1990-2000, 2000-2007
  \item Convert $\hat{\gamma}_{k,t}$ to levels (1970 = 100)
\item $k$ is top five Rotemberg weight industries in 1980, and ``aggregate'' 
	\begin{itemize}
	\item Aggregate: 1980 shares, aggregated using $g_{k,1990-2000}^{\text{high-income}}$
	\end{itemize}

  \end{itemize}
``Pre-period'' prior to 1990

\end{frame}








%%%%%%%%%
% 3571
%%%%%%
%---
\begin{frame}[label=ind3571_gr1980]
\addtocounter{framenumber}{-1}
\textbf{Electronic computers, fixed 1980 industry shares, Rotemberg weight $ 0.183 $}


\begin{center}
		\includegraphics[scale=0.73]{images/share_emp_ind_bhj_3571_pre_trends_index1980.pdf}
\end{center}




\end{frame}



%%%
% 3944
%%%

%---
\begin{frame}[label=ind3944_gr1980]
\textbf{Games and toys, fixed 1980 industry shares, Rotemberg weight $0.138$}


\begin{center}
		\includegraphics[scale=0.73]{images/share_emp_ind_bhj_3944_pre_trends_index1980.pdf}
\end{center}




\end{frame}








%%%
% 3651
%%%%%
%---
\begin{frame}[label=ind3651_gr1980]
\addtocounter{framenumber}{-1}
\textbf{Household audio and video, fixed 1980 industry shares, Rotemberg weight $ 0.085$}


\begin{center}
		\includegraphics[scale=0.73]{images/share_emp_ind_bhj_3651_pre_trends_index1980.pdf}
\end{center}




\end{frame}



%%%
% 3661
%%%

%---
\begin{frame}[label=ind3661_gr1980]
\addtocounter{framenumber}{-1}
\textbf{Telephone apparatus, fixed 1980 industry shares, Rotemberg weight $0.066 $}


\begin{center}
		\includegraphics[scale=0.730]{images/share_emp_ind_bhj_3661_pre_trends_index1980.pdf}
\end{center}




\end{frame}



%%%
% 3577
%%%

%---
\begin{frame}[label=ind3577_gr1980]
\addtocounter{framenumber}{-1}
\textbf{Computer equipment, fixed 1980 industry shares, Rotemberg weight $0.060 $}


\begin{center}
		\includegraphics[scale=0.730]{images/share_emp_ind_bhj_3577_pre_trends_index1980.pdf}
\end{center}




\end{frame}








%%%
% aggregate
%%%

%---
\begin{frame}[label=aggregate_gr1980]
\addtocounter{framenumber}{-1}
\textbf{Aggregate, fixed 1980 industry shares}


\begin{center}
		\includegraphics[scale=0.73]{images/z_pre_trends_index1980.pdf}
\end{center}



\end{frame}



%-----

\begin{frame}[label=overidtest]
\frametitle{Test \#3: Alternative estimators and overidentification tests}

Basic insight: many instruments 
\begin{itemize}
\item Estimators (maximum likelihood): LIML, Hausman, Newey, Woutersen, Chao and Swanson (2012) HFUL (heteroskedasticity-Fuller (1977))
\item Estimators (two-step):  TSLS (problematic), Bartik TSLS, MBTSLS (Anatolyev (2013), and Kolesar et al (2015)) 
\end{itemize} 
Interpretation:
\begin{itemize}
\item Gap between maximum likelihood and two-step estimators is evidence of misspecification 
\end{itemize}
Also, overidentification tests, which provides evidence of misspecification
%
%Note, in Angrist and Krueger (1991), MBTSLS and LIML are very similar



\end{frame}

%-----

\begin{comment}

\begin{frame}[label=overidtest]
\frametitle{Test \#3b: Overidentification tests}

Basic insight: many instruments 
\begin{itemize}
\item Test statistics: for HFUL, Chao, Hausman, Newey, Swanson and Woutersen (2014) test statistic; for TSLS, Sargan; for LIML (Cragg-Donald)
\end{itemize} 
Interpretation:
\begin{itemize}
\item Overid test failure is evidence of misspecification
\end{itemize}
%Note, in Angrist and Krueger (1991), overid tests fail to reject



\end{frame}

\end{comment}

%----
\begin{frame}[fragile,label=sh_results]
\frametitle{Test \#3: Alternative estimators and overidentification}


\small
\begin{table}[htbp]
\label{table:olsiv_bartik}
\begin{center}
  \begin{tabular}{l*{4}{c}}
\toprule
  & \multicolumn{1}{c}{$\Delta$ Emp} & Over ID Test\\
\midrule
OLS  & -0.17 &  \\
    & \textcolor{gray!75}{(0.04)}&         & \\
TSLS (Bartik)  & -0.62 &  \\
              &  \textcolor{gray!75}{(0.11) }&         & \\
TSLS           & -0.22 & 872.69\\
              &  \textcolor{gray!75}{(0.06)}        & [0.00]\\
MBTSLS  & -0.33 & \\
    &  \textcolor{gray!75}{(0.05)}&         & \\
LIML  & -2.07 & 1348.50\\
    &  \textcolor{gray!75}{(3.52)}&          [0.00]\\
HFUL & -1.13 &  1141.08\\
    &  \textcolor{gray!75}{(0.04)}&         [0.00]\\
\midrule
Year and Census Division FE&            Yes\\
Controls        &           Yes\\
Observations    &      1,444\\
\bottomrule
\end{tabular}
\end{center}
\end{table}



\end{frame}


\begin{frame}{Simulated Instruments}
  \begin{wideitemize}
  \item Pivoting now to simulated instruments
  \item Running example -- effect of Medicaid on Y (e.g. child mortality)
  \item Eligibility depends on:
    \begin{itemize}
    \item State-year (legislative changes)
    \item Individual characteristics: Income less disregards (child care and work expenses), adults (if married, unemployed husband can maintain eligibility), number of children, whether first-time mother, etc.  
    \end{itemize} 
    Actual eligibility is clearly endogenous:
    \begin{itemize}
    \item Poorer people are more likely to be eligible, ``worse'' potential outcomes
    \end{itemize}
  \end{wideitemize}
\end{frame}



\begin{frame}
\frametitle{Build an instrument}

\begin{itemize}
\item  Idea: isolate (changes in) eligibility coming from policy changes
  \begin{itemize}
  \item Take a fixed population (say, 1986 CPS)
  \item Compute ``share eligible'' in each state-year 
  \end{itemize} 
\item		Parameterizes (in one-dimension) how generous the policy is
\item Benefits:
  \begin{itemize}
  \item   Intuitive: isolate variation that you think is exogenous  
  \item ``Split-sample'': Can compute simulated eligibility in one sample, and outcomes in another sample
  \end{itemize}
\end{itemize}
\end{frame}


\begin{frame}{Set up}
\begin{wideitemize}
\item $Y_{i}$: outcome (e.g. utilization of health care, whether baby is low birthweight)
\item $D_{i}$: treatment (e.g. Medicaid eligibility)
  \begin{itemize}
  \item Today: $D_{i} \in \{0,1\}$ 
  \end{itemize}
\item $D_{i} \equiv g(S_{i}, X_{i})$: \emph{known} eligibility function $g(\cdot, \cdot)$\\
  \hspace{10pt} $S_{i}$: location  (e.g. state)\\
  \hspace{10pt} $X_{i}$: characteristics driving eligibility (e.g. income, age, \# of children)\\
  \hspace{10pt} $\mathcal{X}$: simplifying case - discrete support of $X_{i}$\\  
\item \emph{$g(s,x)$ can be highly non-linear in $x$}
\end{wideitemize}
\end{frame}

%-----
\begin{frame}
\frametitle{Aronow, Goldsmith-Pinkham, \& Sorkin: treatment effect heterogeneity}

It is natural to imagine that we are not in a constant effects world:
\begin{itemize}
\item In observables: differential effects by income
\item In unobservables: differential effects by latent health and etc.
\end{itemize} 
Two questions with TE heterogeneity:
\begin{itemize}
\item How plausible is it that monotonicity holds?  
\item What is the LATE? (And is the LATE of economic interest?) 
\end{itemize}

\end{frame}

% Some Notation
\begin{frame}{Fundamental problem of causal inference}
  \begin{wideitemize}
  \item $Y_{i}(D_{i})$ denote potential outcomes for individual $i$ 
  \item $Y_{i}(1) - Y_{i}(0)$: individual treatment effect of eligibility
  \item Observe $Y_{i} = D_{i}Y_{i}(1) + (1-D_{i})Y_{i}(0)$ 
  \end{wideitemize}
  Solution:
  \begin{wideitemize}
  \item Assume state \textit{policies} are as good as
    random...
    \begin{itemize} 
    \item Concern that $X_{i}$ is correlated with the
      potential outcomes;
    \item Poor individual more likely to be eligible, but may have worse outcomes ($Y_{i}(D_{i})$);
    \end{itemize}
  \item How to get state policy variation without endogeneity from composition? 
  \end{wideitemize}
\end{frame}

\begin{frame}{Quote - Currie \& Gruber (1996)}
  ``\textbf{Our IV strategy
  ... treats the legislative environment ... as a source of exogenous identification.}''
\end{frame}


\begin{frame}{  The simulated instrument}
 
    \begin{equation*}
      Z_{i} = \sum_{x \in \mathcal{X}}g(S_{i}, x) \underbrace{Pr(X_{i} = x)}_{\text{non-stoch func}}      
    \end{equation*}
    \begin{itemize}
    \item $Pr(X_i=x)$ is, e.g., distribution of characteristics in 1986 CPS
    \end{itemize} 
   TSLS:
    \begin{align*}
      Y_{i} &= \alpha + \beta D_{i} + \epsilon_{i} \\
      D_{i} &= \pi_0 + \pi_{1} Z_{i} + \eta_{i}.
    \end{align*}
    Standard constant effects assumptions:
    \begin{equation*}
      \underbrace{Cov(Z_{i},\epsilon_{i})}_{\text{\textcolor{red}{exclusion}}} = 0, \qquad       \underbrace{Cov(Z_{i},D_{i})}_{\text{\textcolor{blue}{relevance}}} \not= 0
    \end{equation*}
\end{frame}


\begin{frame}{TSLS Estimator and Heterogeneity}
  \begin{wideitemize}
  \item  TSLS estimator:
  \begin{align*}
    \hat{\beta}_{sim} &= \frac{\widehat{Cov}(Z_{i}, Y_{i})}{\widehat{Cov}(Z_{i}, D_{i})}.
  \end{align*}
\item What average TE is $\hat{\beta}_{sim}$? Most researchers appeal
  to monotonicity for LATE:
  \begin{assN}
    Let $\mathcal{Z}$ be the set of values of $Z_{i}$. For all
    $z, z' \in \mathcal{Z}$, either $D_{i}(z) \geq D_{i}(z')$ or
    $D_{i}(z) \leq D_{i}(z')$.
  \end{assN}
\item How plausible is this?
  \end{wideitemize}
\end{frame}


%-----
\begin{frame}
\frametitle{Idea \#1: monotonicity is very likely to be violated}

Recall why we need an instrument: 
\begin{itemize}
\item Eligibility depends on ``high-dimensional'' observables
\item States change many dimensions of eligibility at once 
\item If it were just income, we'd do an RD in income 
\end{itemize}

Hence:
\begin{itemize}
	\item When a state expands eligibility \textbf{on average}, it can make some groups ineligible 
		\begin{itemize}
	\item Law changes represent ``menu of treatments'' 
		\end{itemize}
	\item Violate monotonicity in \textbf{observables}; can document empirical prevalence (!)
	\end{itemize}


Why you probably shouldn't use simulated instruments (in this setting): potentially uninterpretable estimates in presence of TE heterogeneity
\end{frame}

%---
\begin{frame}
\frametitle{Idea \#2: we have a selection on observables design}

\begin{itemize}
\item Basic assumption in the simulated instrument design: state legislative policy changes don't depend on characteristics of the population 
\item \textbf{Strengthen (?) this assumption:} characteristics don't \textbf{respond} to policy changes (i.e., become poorer to become eligible for Medicaid) 
	\begin{itemize}
	\item This form of response affects interpretation of IV estimate 
	\end{itemize}
\item If you buy this (which you might not in some settings!): 
	\begin{itemize}
	\item You can compare outcomes of ``types'' of people (vectors of characteristics) who are eligible in some state-years and not in others
	\end{itemize}
\end{itemize} 

\end{frame}


\begin{frame}{Monotonicity assumption is strong, and testable}
  \begin{wideitemize}
  \item Since $g(\cdot, \cdot)$ is known for all $x$ and $s$, can
    look for $x, x' \in \mathcal{X}$ where $Z_{i}(s) > Z_{i'}(s')$ for
    $x$ and $x'$ but $D_{i}(z) > D_{i}(z')$ for $x$ and
    $D_{i}(z) < D_{i}(z')$
  \item Simple example: states where eligibility is not ordered
    \begin{itemize}
    \item $\mathcal{X} = \{1,2,3,4,5\}$, $Pr(X_{i} = x) = 1/5$
    \item Types 1,2,5 are eligible in $s$, and 1,3,4,5 are eligible in $s'$.
    \item Then, $\forall i$, $Z_{i}(s') > Z_{i}(s)$.
    \item But, for $X_{i} = 3$, $D_{i}(s) < D_{i}(s')$, while for
      $X_{i} = 2$, $D_{i}(s') < D_{i}(s)$
    \end{itemize}
  \item Implication: motonoicity is violated, and $\hat{\beta}_{sim}$
    combines positive and negative eligibility effects.
  \end{wideitemize}
\end{frame}

\begin{frame}{How important are these violations?}
  \begin{wideitemize}
  \item Cohodes, Grossman, Kleiner and Lovenheim (2016) study long-term impact of Medicaid
    expansions on childrens' educational outcomes
  \item Information on eligibility for fixed distribution of 1986 CPS 
  \item Define an eligibility ``type'':
  	\begin{itemize}
	\item A unique pattern of state-year (from 1980-2007) eligibility for a respondent in the 1986 CPS 
	\item 48,036 children in the 1986 CPS, $18,881$ eligibility types 
	\end{itemize}
  \end{wideitemize}
\end{frame}


%-----
\begin{frame}
\frametitle{Quantifying violations of monotonicity}

\begin{itemize}
\item Take state-years where the value of the instrument changes (i.e., state policy becomes more/less generous)
\item Ask for each eligibility type whether its eligibility moved in the direction of the instrument
\item Weight by the eligibility type share of the national population
	\begin{itemize}
	\item In the CPS, small samples at state level
	\end{itemize}
\end{itemize}


\end{frame}


%----
\begin{frame}
\frametitle{Idea \#2.5: What you can sometimes do instead}

With infinite data and outcomes observed in same dataset as eligibility (no split sample): 
	\begin{itemize}
	\item Compute type specific treatment effects 
		\begin{itemize}
		\item Simple comparison of means, no IV! 
		\end{itemize}
	\item Report TE for populations of interest where eligibility varies across state-years
		\begin{itemize}
		\item I.e., if we expanded the income threshold nationally, what would happen...
		\end{itemize} 
	\item Because monotonicity violation is in observables, not hard to fix (!)
	\end{itemize} 
In practice: use the propensity score....
\end{frame}

%-----
\begin{frame}
\frametitle{Caveats}

This is about Currie and Gruber (1996a, b) and not Gruber and Saez (2002):
\begin{itemize}
\item Gruber and Saez is an individual-level instrument so the monotonicity point doesn't apply
\item Gruber and Saez is all about characteristics (income) responding endogenously
\end{itemize} 
Not clear how useful/implementable idea \#2.5 is:
\begin{itemize}
\item Need outcomes and eligibility in same dataset
\end{itemize} 
In progress:
\begin{itemize}
\item A complete empirical example
\item Even if you can't do \#2.5, can bound TE heterogeneity so that still have interpretable estimates
\item Formalism is a bit sloppy: currently just for cross-section, and not panel with two-way F.E.
\end{itemize}

\end{frame}

\begin{frame}{Borusyak and Hull: formalizing this point even further}
  \begin{wideitemize}
  \item Borusyak and Hull (2021) bring this point even further,
    emphasizing this propensity score setup in a wider range of settings
  \item They allow for potential \emph{interference} across designs as well (due to, for example, spatial correlation in things like railroads)
  \item BH approach builds on the BHJ/AKM notion of random shocks
    being combined using non-random exposure
    \begin{itemize}
    \item Key p-score result is the same, but the build out a more
      general (and more efficient) set of structures
    \item Additional, in cases with interference (like railroad
      upgrades), BH show how to ``recenter'' the pscore to account for
      interference
    \end{itemize}
  \end{wideitemize}
\end{frame}



\end{document}
